% Options for packages loaded elsewhere
\PassOptionsToPackage{unicode}{hyperref}
\PassOptionsToPackage{hyphens}{url}
\documentclass[
]{article}
\usepackage{xcolor}
\usepackage{amsmath,amssymb}
\setcounter{secnumdepth}{-\maxdimen} % remove section numbering
\usepackage{iftex}
\ifPDFTeX
  \usepackage[T1]{fontenc}
  \usepackage[utf8]{inputenc}
  \usepackage{textcomp} % provide euro and other symbols
\else % if luatex or xetex
  \usepackage{unicode-math} % this also loads fontspec
  \defaultfontfeatures{Scale=MatchLowercase}
  \defaultfontfeatures[\rmfamily]{Ligatures=TeX,Scale=1}
\fi
\usepackage{lmodern}
\ifPDFTeX\else
  % xetex/luatex font selection
\fi
% Use upquote if available, for straight quotes in verbatim environments
\IfFileExists{upquote.sty}{\usepackage{upquote}}{}
\IfFileExists{microtype.sty}{% use microtype if available
  \usepackage[]{microtype}
  \UseMicrotypeSet[protrusion]{basicmath} % disable protrusion for tt fonts
}{}
\makeatletter
\@ifundefined{KOMAClassName}{% if non-KOMA class
  \IfFileExists{parskip.sty}{%
    \usepackage{parskip}
  }{% else
    \setlength{\parindent}{0pt}
    \setlength{\parskip}{6pt plus 2pt minus 1pt}}
}{% if KOMA class
  \KOMAoptions{parskip=half}}
\makeatother
\usepackage{longtable,booktabs,array}
\newcounter{none} % for unnumbered tables
\usepackage{multirow}
\usepackage{calc} % for calculating minipage widths
% Correct order of tables after \paragraph or \subparagraph
\usepackage{etoolbox}
\makeatletter
\patchcmd\longtable{\par}{\if@noskipsec\mbox{}\fi\par}{}{}
\makeatother
% Allow footnotes in longtable head/foot
\IfFileExists{footnotehyper.sty}{\usepackage{footnotehyper}}{\usepackage{footnote}}
\makesavenoteenv{longtable}
\usepackage{graphicx}
\makeatletter
\newsavebox\pandoc@box
\newcommand*\pandocbounded[1]{% scales image to fit in text height/width
  \sbox\pandoc@box{#1}%
  \Gscale@div\@tempa{\textheight}{\dimexpr\ht\pandoc@box+\dp\pandoc@box\relax}%
  \Gscale@div\@tempb{\linewidth}{\wd\pandoc@box}%
  \ifdim\@tempb\p@<\@tempa\p@\let\@tempa\@tempb\fi% select the smaller of both
  \ifdim\@tempa\p@<\p@\scalebox{\@tempa}{\usebox\pandoc@box}%
  \else\usebox{\pandoc@box}%
  \fi%
}
% Set default figure placement to htbp
\def\fps@figure{htbp}
\makeatother
\usepackage{svg}
\setlength{\emergencystretch}{3em} % prevent overfull lines
\providecommand{\tightlist}{%
  \setlength{\itemsep}{0pt}\setlength{\parskip}{0pt}}
\usepackage{bookmark}
\IfFileExists{xurl.sty}{\usepackage{xurl}}{} % add URL line breaks if available
\urlstyle{same}
\hypersetup{
  hidelinks,
  pdfcreator={LaTeX via pandoc}}

\author{}
\date{}

\begin{document}

\textsc{Comparative Reproductive Biology And Impacts Of Selfing In Two
Epiphytic Bromeliaceae}

\textsc{Joshua M. Felton\textsuperscript{*}}

School of Integrative Plant Science, Section of Plant Biology and the L.
H. Bailey Hortorium, Cornell University, Ithaca, New York, USA

\textsc{M. Shane Heschel}

Department of Organismal Biology \& Ecology, Colorado College, 14 E.
Cache la Poudre Street, Colorado Springs, Colorado, USA 80903

\textsc{Erin N. Bodine}

Department of Mathematics and Statistics, Rhodes College, 2000 N.
Parkway Ave, Memphis, Tennessee, USA 39112

\textsc{Rachel S. Jabaily}

Department of Organismal Biology \& Ecology, Colorado College, 14 E.
Cache la Poudre Street, Colorado Springs, Colorado, USA 80903

*Corresponding author: jmf425@cornell.edu

\textsc{Abstract.} The breeding systems of most Bromeliaceae have not
been tested directly. We characterized the floral biology and breeding
systems of two epiphytic, tank-forming bromeliads endemic to the
Atlantic Forest of Brazil, \emph{Vriesea rafaelii} (Tillandsioideae) and
\emph{Billbergia brasiliensis} (Bromelioideae), using controlled self-
and cross-pollinations in a greenhouse. Both species were
self-compatible. Based on seed-set, the self-compatibility index was
0.59 in \emph{V. rafaelii} and 0.87 in \emph{B. brasiliensis}. In
\emph{B. brasiliensis}, selfed flowers produced lighter fruits and
lighter seeds than outcrossed flowers, whereas in \emph{V. rafaelii}
fruit and seed mass did not differ between self- and cross-pollinated
flowers. A subset of \emph{V. rafaelii} flowers showed a spatial
separation of the anthers and stigma (herkogamy) that prevented
autonomous selfing. These results add two species to the small set of
bromeliads with experimentally characterized breeding systems and
highlight differences in reproductive output with selfing.

\emph{Keywords:} \emph{Billbergia brasiliensis}, breeding system,
controlled pollination, herkogamy, self-compatibility, \emph{Vriesea
rafaelii}

\textbf{\textsc{Introduction}}

A plant\textquotesingle s breeding system, whether its flowers are
self-compatible or self-incompatible, is

a fundamental part of its life history, determining seed set and
pollination ecology. Many plants are self compatible and can
self-pollinate (Igic \& Busch 2013), which can be advantageous when
mates or pollinators are scarce because it provides reproductive
assurance. However, selfing can be costly due to reductions in
heterozygosity, which can result in the expression of deleterious
recessive alleles. This often produces inbreeding depression, in which
selfed offspring set fewer or lighter fruits and fewer or less viable
seeds than outcrossed offspring (Hedrick \& Kalinowski 2000; Angeloni et
al. 2011). Determining whether a species can self, and whether selfing
reduces its reproductive output, contextualizes the evolutionary trends
of breeding systems and builds natural history knowledge. The
Neotropical family Bromeliaceae is large and varied in floral form, and
its breeding systems span the full range from obligate selfing to
obligate outcrossing, either through self-incompatibility (SI)
mechanisms or through dioecy (reviewed in Cascante-Marín \&
Núñez-Hidalgo 2023). Self-compatibility is common, occurring in nearly
all genera of subfamily Tillandsioideae and in only a few genera of
Bromelioideae, the two subfamilies that are the focus of this study.
Some species also vary in flower size and in the arrangement of their
reproductive organs within a single inflorescence, with certain flowers
favoring outcrossing through a spatial separation of the anthers and
stigma (herkogamy). For most bromeliads, however, the breeding system
has never been directly tested (Cascante-Marín \& Núñez-Hidalgo 2023).
Here we characterize and compare the reproductive biology of two
epiphytic, tank-forming bromeliads endemic to the Atlantic Forest of
Brazil. \emph{Vriesea rafaelii} Leme (subfamily Tillandsioideae)
produces an erect spike inflorescence and dry, capsular fruits, while
\emph{Billbergia brasiliensis} L.B.Sm. (subfamily Bromelioideae)
produces a pendant spike inflorescence and fleshy fruits. We (a)
describe the reproductive biology of each species, including variation
among individual flowers, and (b) test whether the biomass a plant
invests in fruits and seeds differs between selfed and outcrossed
flowers.

\textbf{\textsc{Methods}}

The taxa were purchased from Tropiflora (Sarasota, FL, USA) and grown
from seed collected from a single parental plant of each species. Plant
care and earlier developmental data are described in (Jabaily et al.
2021). Seventeen~\emph{V. rafaelii} and 14~\emph{B.
brasiliensis}~individuals with visible developing inflorescences were
included in the study.~Within each species, individuals were assigned to
a selfing or an outcrossing treatment (\emph{V. rafaelii}: 9 self, 8
outcross; \emph{B. brasiliensis}: 7 self, 7 outcross). Flowers left
unmanipulated were scored as autonomously selfed.

\emph{V. rafaelii}~is protogynous: the stigma becomes receptive shortly
after the petals open, and the anthers dehisce about three hours later.
We monitored receptivity with a hand lens, noting receptivity once the
stigma was visibly wet with secretions from its surface papillae. For
the outcrossing treatment, pollen was applied to receptive stigmas in
the evening, and the anthers were removed before dehiscence to prevent
self-pollination. When only one flower was available on an individual (N
= 9 flowers), the excised anthers were stored at -20°C and used later as
donor pollen, as frozen bromeliad pollen remains viable for crossing (De
Souza et al. 2015).

\emph{Billbergia brasiliensis}~anthers and stigma mature together and
remain in proximity during anthesis. Pollinations were done in late
morning to coincide with anther dehiscence. For the outcrossing
treatment, flowers were opened by hand before the petals curled back,
the anthers were removed to prevent accidental self-pollination, and the
flower was enclosed in a fine mesh organza bag, which helped ensure no
accidental crossing occurred from pollen from other plants, until the
controlled cross was performed.

Fruits of \emph{B. brasiliensis} and \emph{V. rafaelii} require roughly
seven and eighteen months, respectively, to reach full maturity (a color
change from green to yellow in \emph{B. brasiliensis}, and the onset of
capsule dehiscence in \emph{V. rafaelii}). Because of time constraints
on data collection, inflorescences were harvested before full ripeness,
after about four months for~\emph{B. brasiliensis}~and six months
for~\emph{V. rafaelii}. We considered a flower aborted if it failed to
produce any seed (seed mass = 0 g). Harvested inflorescences were dried
in paper bags for five days at 60°C, and the rachis, scape bracts, and
floral bracts were weighed separately from the fruits. Each dried fruit
was dissected to separate seeds and their associated structures from
fruit tissue, and seeds and fruit were each weighed on an analytical
balance, with seed mass recorded as the total mass of all seeds per
fruit rather than the mass of an individual seed. To quantify each
species' capacity for self-fertilization, we calculated a
Self-Compatibility Index (Zapata \& Arroyo 1978).

\[\begin{array}{r}
Self\ Compatibility\ Index = P_{a}\ /{\ P}_{x}\ \#(1)
\end{array}\]

where \(P_{a}\) is the proportion of selfed flowers that set fruit and
\(P_{x}\) is the proportion of cross-pollinated flowers that set fruit.
The index ranges from 0 (completely self-incompatible) to 1 (fully
self-compatible). The index assumes Pa ≤ Px; values above 1 would
reflect experimental error rather than genuine self-compatibility. To
test whether cross-type (selfed vs. outcrossed) influenced fruit mass or
seed mass, we used a nested mixed-model ANOVA with species, cross-type,
and their interaction as fixed effects and individual plant identity
nested within species as a random effect. Analyses were performed in JMP
(SAS Institute).

\textbf{\textsc{Results}}

\emph{Natural history. Billbergia brasiliensis} had a spike
inflorescence that was pendant and white from farinose hairs
(\textsc{Figure 1a}). Stigma receptivity occurred at the same time of
anther dehiscence which had occurred by the time the corolla has
spiraled into a coil (\textsc{Figure 1b-c}). Flowers did not exhibit any
form of herkogamy (\textsc{Figure 1c}). The typical showy pink bracts
varied in pigment between individuals. Multiple individuals had bracts
that were brown and dried at anthesis (\textsc{Figure 1d-e}). Multiple
individual scapes did not elongate enough to become pendant and instead
flowered when erect, primarily within the tubular vegetative rosette.
Flowers on these scapes were difficult to observe and manipulate. Fruits
changed pigment from green to yellow once ripe (at least six months post
anthesis) and were subrotund with ribs. Seeds were coated in a
mucilaginous substance (\textsc{Figure 1f}).

\emph{Vriesea rafaelii} had an erect spike inflorescence with dark green
scape and floral bracts (\textsc{Figure 1g}). Stigma receptivity began
shortly after petals had opened, with anthers dehiscing approximately
three hours after petal opening (\textsc{Figure 1i}). Anthers were
arranged ventrally in the flower, touching the stigma (\textsc{Figure
1h}). Most individuals had anthers at the same length as the stigma.
Mismatch in anther and pistil length was noted particularly in
individuals (\emph{N} = 10 individuals) who began their development
later in the study population. The most common mismatch observed was the
stamen being shorter than the pistil length, so much so that selfing did
not occur (\emph{N} = 29 flowers) (\textsc{Figure 1j}). Petal
appendages, which are a single flap of tissue between the staminal
filament and the petal at the base of the corolla (Brown \& Terry 1992),
occasionally were long enough to cover the stigma (\emph{N} = 4
flowers), leading to no fruit set from selfing (\textsc{Figure 1k}).
Capsule development took about seventeen months from flowering
(\textsc{Figure 1l}).

\emph{Breeding system.} All hand-pollinated flowers formed fruits
regardless of cross type. However, successful seed production was lower
following selfing than outcrossing (81\% vs. 93\% in \emph{B.
brasiliensis} and 34\% vs. 58\% in \emph{V. rafaelii}), indicating that
the primary reproductive cost of selfing occurred after fruit initiation
(Table 1). The effect of cross type on capsule mass and on seed mass did
not differ significantly between the two species (capsule mass, F =
1.29, df = 1, P = 0.257; seed mass, F = 1.19, df = 1, P = 0.276), so we
report each species separately. In~\emph{Billbergia brasiliensis},
outcrossed flowers produced heavier fruits (F = 4.08, df = 1, P = 0.045)
and heavier seeds (F = 7.23, df = 1, P = 0.008) than selfed flowers
(\textsc{Figure 2a}). In~\emph{Vriesea rafaelii}, capsule mass was
marginally lower in outcrossed than selfed flowers but did not differ
significantly between cross types (F = 2.48, df = 1, P = 0.116;
\textsc{Table S1}), and seed mass did not differ between cross types
when flowers that set no seed are included (F = 0.85, df = 1, P = 0.356;
\textsc{Table S2}; \textsc{Figure 2b}). Both species were
self-compatible, with a self-compatibility index of 0.87 in~\emph{B.
brasiliensis} and 0.59 in~\emph{V. rafaelii}~(\textsc{Table 2}).

\textbf{\textsc{Discussion}}

Both species were self-compatible, but selfing was costlier in more ways
in \emph{B. brasiliensis} than in \emph{V. rafaelii}. In~\emph{B.
brasiliensis}, selfed flowers produced lighter fruits and lighter seeds
than outcrossed flowers and set seed somewhat less often, indicating
early-acting inbreeding depression in both the formation and the
provisioning of seed. Reduced fruit set and smaller seeds following
self-pollination have likewise been reported in field studies
of~\emph{Canistrum aurantiacum} and~\emph{Encholirium heloisae}
(Christianini et al. 2013)(Siqueira Filho \& Machado 2001),
placing~\emph{B. brasiliensis}~among the bromeliads in which selfing
lowers early reproductive output. In \emph{V. rafaelii}, the cost of
selfing was confined to a single step. Fruit and seed mass did not
differ between cross types, but selfed flowers were far less likely to
set any seed at all than outcrossed flowers (34\% versus 58\%). Among
the flowers that did set seed, the seed crop was no smaller after
selfing than after outcrossing. The lower reproductive output of selfed
\emph{V. rafaelii} therefore came largely from flowers that failed to
set seed, which might be indicative of self-compatibility issues and
higher abortion rates due to inbreeding. Mean seed mass alone therefore
understated inbreeding depression in~\emph{V. rafaelii}, because much of
the effect was carried by the many selfed flowers that set no seed. A
subset of \emph{V. rafaelii} flowers was also herkogamous, with the
anthers held away from the stigma or a petal appendage covering it, and
these flowers formed fruits and set seed only when outcrossed, although
this was observed in only three flowers. Because fruit and seed mass
index an early stage of offspring performance, these measures bound
rather than resolve the total strength of inbreeding depression in
either species; resolving it will require following selfed and
outcrossed progeny to germination and later life-history stages (Riginos
et al. 2007) as well as estimating inbreeding directly with molecular
markers (Matos Paggi et al. 2015). These results add two previously
uncharacterized species to the small set of bromeliads with
experimentally tested breeding systems (Cascante-Marín \& Núñez-Hidalgo
2023). Because both species can self yet incur a cost at the seed stage,
selfing and inbreeding may impact population dynamics.~ However, the
degree to which selfing contributes to population persistence and
recruitment in the wild remains an open question that greenhouse crosses
can only begin to answer.

Both species fell well within the self-compatible range by an index of
self-incompatibility, though~\emph{V. rafaelii}~was the less
self-compatible of the two (\textsc{Table 2}). Our finding runs counter
to the general expectation for the family: self-compatibility is thought
to be widespread in Tillandsioideae and comparatively rare in
Bromelioideae (Cascante-Marín \& Núñez-Hidalgo 2023), yet here the
bromelioid \emph{B. brasiliensis} was more self-compatible than the
tillandsioid \emph{V. rafaelii}. This suggests that breeding system in
Bromeliaceae can vary even within lineages generally considered
self-incompatible, and underscores how little direct testing exists
outside a handful of species. The within-plant variation in herkogamy
in~\emph{V. rafaelii}~also shows that the capacity to self can vary
among flowers on a single individual, which is relevant to how selfing
rates are estimated from whole-plant or whole-population data (Whitehead
et al. 2018). Self-compatibility promotes reproductive assurance in the
absence of mates or pollinators, but selfing carried a measurable cost
in both species, expressed as reduced seed set in \emph{V. rafaelii} and
as both reduced seed set and lighter seeds in \emph{B. brasiliensis}.
Therefore, the capacity to self does not translate directly into
equivalent reproductive output, indicating inbreeding depression and
potential impacts on population dynamics. Future work should investigate
this further on three fronts: following selfed and outcrossed progeny
beyond seed to germination and establishment, where later-acting
inbreeding depression may yet emerge; estimating realized outcrossing
rates in natural populations with molecular markers; and documenting the
floral visitors and pollination behavior that determine how often
outcrossing occurs at these sites. For two narrowly distributed
endemics, resolving how much of their reproduction is selfed, and at
what cost, would do much to clarify how secure their mating systems
leave them.

\textbf{Acknowledgements}

The authors would like to thank Ali Keller for plant care assistance,
Dr. Brad Oberle and Dr. Brian Sidoti for their insights and comments
during the project, and the Colorado College Keller family fund and the
Hevey fund for financial support.

\textbf{Authors Contribution Statement}

JMF, MSH, RSJ conceptualized the project; JMF collected the data; JMF,
MSH, ENB analyzed the data; JMF, RSJ wrote the original draft; all
authors reviewed and edited the final draft

\textbf{ORCID}

JMF: 0000-0002-2882-0655; RSJ: 0000-0002-4218-3533; ENB:
0000-0001-7885-1233; MSH: 0000-0003-4455-0920

\textbf{References}

Angeloni, F., Ouborg, N. J., and Leimu, R. 2011. Meta-analysis on the
association of population size and life history with inbreeding
depression in plants. Biological Conservation 144: 35--43.
https://doi.org/10.1016/j.biocon.2010.08.016

Brown, G. K. and Terry, R. G. 1992. Petal Appendages in Bromeliaceae.
American Journal of Botany 79: 1051--1071.
https://doi.org/10.2307/2444915

Cascante-Marín, A. and Núñez-Hidalgo, S. 2023. A Review of Breeding
Systems in the Pineapple Family (Bromeliaceae, Poales). The Botanical
Review 89: 308--329. https://doi.org/10.1007/s12229-023-09290-0

Christianini, A. V., Forzza, R. C., and Buzato, S. 2013. Divergence on
floral traits and vertebrate pollinators of two endemic
\emph{Encholirium} bromeliads. Plant Biology 15: 360--368.
https://doi.org/10.1111/j.1438-8677.2012.00649.x

De Souza, E. H., Souza, F. V. D., Rossi, M. L., Brancalleão, N., Da
Silva Ledo, C. A., and Martinelli, A. P. 2015. Viability, storage and
ultrastructure analysis of Aechmea bicolor (Bromeliaceae) pollen grains,
an endemic species to the Atlantic forest. Euphytica 204: 13--28.
https://doi.org/10.1007/s10681-014-1273-3

Hedrick, P. W. and Kalinowski, S. T. 2000. Inbreeding Depression in
Conservation Biology. Annual Review of Ecology and Systematics 31:
139--162. https://doi.org/10.1146/annurev.ecolsys.31.1.139

Igic, B. and Busch, J. W. 2013. Is self‐fertilization an evolutionary
dead end? New Phytologist 198: 386--397.
https://doi.org/10.1111/nph.12182

Jabaily, R. S., Oberle, B., Fetterly, E. W., Heschel, M. S., Sidoti, B.
J., and Bodine, E. N. 2021. Refining Iteroparity with Comparative
Morphometric Data in Bromeliaceae. International Journal of Plant
Sciences 182: 577--590. https://doi.org/10.1086/715484

Matos Paggi, G., Palma-Silva, C., Bodanese-Zanettini, M. H., Lexer, C.,
and Bered, F. 2015. Limited pollen flow and high selfing rates toward
geographic range limit in an Atlantic forest bromeliad. Flora -
Morphology, Distribution, Functional Ecology of Plants 211: 1--10.
https://doi.org/10.1016/j.flora.2015.01.001

Riginos, C., Heschel, M. S., and Schmitt, J. 2007. Maternal effects of
drought stress and inbreeding in \emph{Impatiens capensis}
(Balsaminaceae). American Journal of Botany 94: 1984--1991.
https://doi.org/10.3732/ajb.94.12.1984

Siqueira Filho, J. A. D. and Machado, I. C. S. 2001. Biologia
reprodutiva de Canistrum aurantiacum E. Morren (Bromeliaceae) em
remanescente da Floresta Atlântica, Nordeste do Brasil. Acta Botanica
Brasilica 15: 427--443. https://doi.org/10.1590/S0102-33062001000300011

Whitehead, M. R., Lanfear, R., Mitchell, R. J., and Karron, J. D. 2018.
Plant Mating Systems Often Vary Widely Among Populations. Frontiers in
Ecology and Evolution 6: 38. https://doi.org/10.3389/fevo.2018.00038

Zapata, T. R. and Arroyo, M. T. K. 1978. Plant Reproductive Ecology of a
Secondary Deciduous Tropical Forest in Venezuela. Biotropica 10: 221.
https://doi.org/10.2307/2387907

\textsc{\hfill\break
Table 1.} Breeding-system descriptors for hand-pollinated flowers of
\emph{Billbergia brasiliensis} and \emph{Vriesea rafaelii}. Fruit and
seed mass are mean ± SD. Fruit set is the proportion of flowers forming
a capsule; seed set is the proportion producing seed. A flower was
scored as aborted if it failed to set any seed (seed mass = 0 g). ``All
flowers'' rows include aborted flowers; ``non-aborted flowers'' rows
exclude them and correspond to the values plotted in Fig. 2 and analyzed
in Tables S1--S2.

{\def\LTcaptype{none} % do not increment counter
\begin{longtable}[]{@{}
  >{\raggedright\arraybackslash}p{(\linewidth - 8\tabcolsep) * \real{0.2650}}
  >{\raggedright\arraybackslash}p{(\linewidth - 8\tabcolsep) * \real{0.1838}}
  >{\raggedright\arraybackslash}p{(\linewidth - 8\tabcolsep) * \real{0.1838}}
  >{\raggedright\arraybackslash}p{(\linewidth - 8\tabcolsep) * \real{0.1838}}
  >{\raggedright\arraybackslash}p{(\linewidth - 8\tabcolsep) * \real{0.1838}}@{}}
\toprule\noalign{}
\multirow{2}{=}{\begin{minipage}[b]{\linewidth}\raggedright
Measure
\end{minipage}} &
\multicolumn{2}{>{\raggedright\arraybackslash}p{(\linewidth - 8\tabcolsep) * \real{0.3675} + 2\tabcolsep}}{%
\begin{minipage}[b]{\linewidth}\raggedright
\emph{Billbergia brasiliensis}
\end{minipage}} &
\multicolumn{2}{>{\raggedright\arraybackslash}p{(\linewidth - 8\tabcolsep) * \real{0.3675} + 2\tabcolsep}@{}}{%
\begin{minipage}[b]{\linewidth}\raggedright
\emph{Vriesea rafaelii}
\end{minipage}} \\
& \begin{minipage}[b]{\linewidth}\raggedright
Outcross
\end{minipage} & \begin{minipage}[b]{\linewidth}\raggedright
Self
\end{minipage} & \begin{minipage}[b]{\linewidth}\raggedright
Outcross
\end{minipage} & \begin{minipage}[b]{\linewidth}\raggedright
Self
\end{minipage} \\
\midrule\noalign{}
\endhead
\bottomrule\noalign{}
\endlastfoot
Flowers, n & 55 & 190 & 43 & 253 \\
Plants, n & 7 & 7 & 8 & 9 \\
Fruit set, \% & 100 & 100 & 100 & 100 \\
Seed set, \% & 93 & 81 & 58 & 34 \\
Fruit mass, all flowers g & 0.146 ± 0.034 & 0.130 ± 0.048 & 0.514 ±
0.250 & 0.406 ± 0.301 \\
Seed mass, all flowers g & 0.168 ± 0.091 & 0.137 ± 0.112 & 0.110 ± 0.119
& 0.082 ± 0.141 \\
Fruit mass, non-aborted flowers, g & 0.153 ± 0.026 & 0.147 ± 0.034 &
0.706 ± 0.130 & 0.792 ± 0.152 \\
Seed mass, non-aborted flowers g & 0.182 ± 0.081 & 0.170 ± 0.100 & 0.189
± 0.096 & 0.239 ± 0.142 \\
\end{longtable}
}

\textsc{Table 2.} Indices of self-compatibility and inbreeding
depression for \emph{B. brasiliensis} and \emph{V. rafaelii}. SC-index
(Eq. 1) calculated as selfed set ÷ outcrossed set; values near 1
indicate self-compatibility and near 0 self-incompatibility, with the
conventional threshold for self-compatibility at approximately 0.2
(Zapata \& Arroyo 1978).

{\def\LTcaptype{none} % do not increment counter
\begin{longtable}[]{@{}
  >{\raggedright\arraybackslash}p{(\linewidth - 4\tabcolsep) * \real{0.3590}}
  >{\raggedright\arraybackslash}p{(\linewidth - 4\tabcolsep) * \real{0.3205}}
  >{\raggedright\arraybackslash}p{(\linewidth - 4\tabcolsep) * \real{0.3205}}@{}}
\toprule\noalign{}
\begin{minipage}[b]{\linewidth}\raggedright
Index
\end{minipage} & \begin{minipage}[b]{\linewidth}\raggedright
\emph{Billbergia brasiliensis}
\end{minipage} & \begin{minipage}[b]{\linewidth}\raggedright
\emph{Vriesea rafaelii}
\end{minipage} \\
\midrule\noalign{}
\endhead
\bottomrule\noalign{}
\endlastfoot
SC-index (fruit set) & 1.00 & 1.00 \\
SC-index (seed set) & 0.87 & 0.59 \\
\end{longtable}
}

\includegraphics[width=6.5in,height=2.6in]{./media/image1.png}

\includesvg[width=6.5in,height=3.38264in]{./media/image3.svg}

\textsc{Figure 1.} Reproductive morphology of \emph{Billbergia
brasiliensis} (a--f) and \emph{Vriesea rafaelii} (g--l). \textbf{A.}
Pendant spike inflorescence with showy pink floral bracts, the axis
white with farinose hairs. \textbf{B.} Single flower at anthesis (ca. 65
mm long; Smith \& Downs 1979) with the corolla coiled, the stigma
receptive as the anthers dehisce. \textbf{C.} Open flower (sepals ca. 10
mm long; Smith \& Downs 1979) showing the anthers and stigma in contact,
with no spatial separation of the two (no herkogamy); stamens are
shorter than the petals in this species, consistent with the observed
contact. \textbf{D. E.} Individuals whose scapes did not elongate, so
the inflorescence flowered erect within the leaf rosette, with floral
bracts brown and dried at anthesis (arrows). \textbf{F.} Ripe fruits,
which turn from green to yellow and are subrotund and ribbed, and seeds
coated in mucilage. \textbf{G.} Erect spike inflorescence with dark
green scape and floral bracts. \textbf{H.} Flower in face view (petals
ca. 48 × 30 mm; Leme 1999), the anthers held against the stigma.
\textbf{I.} Flower in lateral view (flower ca. 65 mm long; Leme 1999);
anthers dehisce about three hours after the petals open. \textbf{J.}
Flower in which the stamens are shorter than the pistil (arrow), so
self-pollination did not occur. \textbf{K.} Flower in which a petal
appendage (cymbiform, ca. 12--14 × 6 mm; Leme 1999) extends across the
stigma (arrow), preventing self-pollination. \textbf{L.} Dehisced
capsules and plumose seeds at maturity, about seventeen months after
flowering.

\textsc{Figure 2.} Dry mass of fruits and seeds from selfed and
outcrossed flowers of \textbf{A.} \emph{Billbergia brasiliensis} and
\textbf{B.} \emph{Vriesea rafaelii}. In each panel the left pair of
boxes is fruit mass, and the right pair is total seed mass including
appendages; selfed flowers are shown in grey and outcrossed flowers in
green. The number of flowers (n) is given above each box. P values are
for the effect of cross type within each species, from the mixed-model
ANOVAs in \textsc{Table S1} (fruit mass) and \textsc{Table S2} (seed
mass). The points show the fruit mass of aborted flowers (\emph{B.
brasiliensis}: 37 selfed, 4 outcrossed; \emph{V. rafaelii}: 166 selfed,
18 outcrossed).

\end{document}
